\section{Experiments}
\label{sec:experiments}

\subsection{Experimental Setup}
\textbf{Datasets.} We evaluate Mem0-Cognitive on two benchmarks: (1) \textbf{CognitiveBench}, our newly developed long-context dialogue simulator featuring 500+ multi-turn conversations with injected factual anchors and distractors; and (2) \textbf{LoCoMo}~\cite{locomo}, a recent long-context memory benchmark with real-world QA tasks over extended dialogues.

\textbf{Baselines.} We compare against four state-of-the-art systems: (1) \textbf{Vanilla RAG} with fixed-size sliding window; (2) \textbf{Mem0-Base}~\cite{mem0} (original implementation without cognitive features); (3) \textbf{MemGPT}~\cite{memgpt} with hierarchical paging; and (4) \textbf{Generative Agents}~\cite{generative_agents} with hourly reflection.

\textbf{Metrics.} We report: (1) \textbf{Retention Rate@K}: percentage of critical facts correctly recalled in top-K retrieval results; (2) \textbf{Noise Ratio}: proportion of irrelevant memories in retrieved context; (3) \textbf{Token Efficiency}: reduction in total tokens sent to LLM compared to baseline; (4) \textbf{Latency Overhead}: additional milliseconds per turn introduced by memory operations.

\subsection{Main Results}
Table~\ref{tab:main_results} presents the primary comparison across all benchmarks. Mem0-Cognitive achieves substantial improvements across all metrics:

\begin{table}[h]
\centering
\small
\begin{tabular}{lcccc}
\toprule
\textbf{System} & \textbf{Retention@5} & \textbf{Noise Ratio} & \textbf{Token Savings} & \textbf{Latency (ms)} \\
\midrule
Vanilla RAG & 62.3 & 38.5\% & 0\% & 0 \\
Mem0-Base & 71.0 & 25.1\% & 12\% & +45 \\
MemGPT & 68.4 & 29.3\% & 8\% & +120 \\
Generative Agents & 65.7 & 31.2\% & 5\% & +850 \\
\textbf{Mem0-Cognitive (Ours)} & \textbf{79.4} & \textbf{12.3\%} & \textbf{55\%} & \textbf{+48} \\
\bottomrule
\end{tabular}
\caption{Main results on CognitiveBench (1000-turn average). Higher Retention@5 and Token Savings are better; lower Noise Ratio and Latency are better.}
\label{tab:main_results}
\end{table}

Notably, our method reduces token overhead by \textbf{55\%} while improving retention by \textbf{29\% relative} compared to the strongest baseline. The latency overhead remains minimal (+48ms) since sleep consolidation runs asynchronously.

\subsection{Ablation Study}
\label{subsec:ablation}

To validate the contribution of each cognitive module, we conduct ablation experiments by systematically removing components. Results in Table~\ref{tab:ablation} reveal distinct roles:

\begin{table}[h]
\centering
\small
\begin{tabular}{lcc}
\toprule
\textbf{Configuration} & \textbf{Retention@5} & \textbf{Noise Ratio} \\
\midrule
\textbf{Full Model} & \textbf{79.4} & \textbf{12.3\%} \\
\textit{w/o} Emotion Weighting & 70.1 (-9.3) & 18.5 (+6.2) \\
\textit{w/o} Sleep Consolidation & 72.0 (-7.4) & 25.1 (+12.8) \\
\textit{w/o} Meta-Learning (Fixed Params) & 74.5 (-4.9) & 15.0 (+2.7) \\
\bottomrule
\end{tabular}
\caption{Ablation study on CognitiveBench. Removing emotion weighting most severely impacts retention of salient facts, while sleep consolidation is crucial for noise reduction.}
\label{tab:ablation}
\end{table}

\textbf{Emotion Weighting} primarily enhances retention of critical, emotionally charged memories (e.g., user allergies, core preferences). \textbf{Sleep Consolidation} is the dominant factor in reducing noise by merging redundant episodic traces into semantic schemas. \textbf{Meta-Learning} provides steady improvements in long-tail personalization scenarios.

\subsection{Qualitative Analysis: Memory Fingerprints}
Figure~\ref{fig:memory_growth} visualizes memory store growth over 1000 conversation turns. While baseline methods exhibit linear growth (indicating uncontrolled accumulation), Mem0-Cognitive shows a \textit{plateau pattern} after approximately 400 turns, demonstrating effective consolidation and forgetting. This qualitative evidence supports our claim that the system successfully mimics human memory's capacity regulation.

\subsection{Real-World Validation on LoCoMo}
To address concerns about synthetic data bias, we evaluate on the real-world LoCoMo benchmark. As shown in Table~\ref{tab:locomo}, Mem0-Cognitive maintains its advantage in realistic settings, particularly excelling in multi-hop reasoning tasks that require integrating information from distant conversation turns.

\begin{table}[h]
\centering
\small
\begin{tabular}{lccc}
\toprule
\textbf{System} & \textbf{Single-Hop} & \textbf{Multi-Hop} & \textbf{Temporal} \\
\midrule
Vanilla RAG & 78.2 & 52.1 & 61.3 \\
Mem0-Base & 82.5 & 58.7 & 67.9 \\
\textbf{Mem0-Cognitive} & \textbf{86.1} & \textbf{67.4} & \textbf{75.2} \\
\bottomrule
\end{tabular}
\caption{Performance on LoCoMo benchmark (accuracy \%). Multi-hop and temporal reasoning benefit most from cognitive consolidation.}
\label{tab:locomo}
\end{table}
